\section{Smooth coarea metrics and a dominated base}
\label{coa:analytic-tools}

Integration along a fibration turns an anticanonical metric upstairs into
an adjoint metric on its base.  In this section the coefficient metric is
smooth.  This regularity gives a direct curvature calculation and, by
coarea, an integrability exponent strictly greater than one.  Together
these facts produce a boundary correction supported only on base divisors
whose inverse images are exceptional over the original variety.

We first establish the analytic statement, keeping curvature and
integrability separate.  We then construct a maximal base dominated by
the anticanonical class and verify the metric hypotheses on that base.
The controlled-boundary criterion converts the resulting numerical
pseudoeffectivity into rational linear effectivity.

\subsection{The primitive representative and horizontal strictness}

\begin{lemma}\label{coa:curvature}
Let $f:\mathcal Z\to\mathbb D$ be a proper holomorphic submersion from a
K\"ahler manifold, with fibers of complex dimension $r>0$.  Let
$(\mathcal L,h)$ be a holomorphic line bundle with a smooth Hermitian metric
of semipositive curvature.  Suppose
$\mathcal E=f_*(K_{\mathcal Z/\mathbb D}+\mathcal L)$ is locally free and
commutes with base change.  Give $\mathcal E$ its fiberwise $L^2$ metric.
Write $D'$ for the $(1,0)$ part of the Chern connection of this
direct-image metric.  For a local holomorphic section $v$ with
$D'v(0)=0$, there is a smooth
absolute $\mathcal L$-valued $(r,0)$-form $w$ representing $v$ such that
\begin{equation}\label{coa:curvature-bound}
 -i\partial\bar\partial\|v\|^2\big|_0
 \ \geq\
 f_*\bigl(c_r w\wedge\overline w\,h\wedge i\Theta_{\mathcal L,h}\bigr)
 \big|_0,\qquad c_r=i^{r^2}.
\end{equation}
In particular, if $i\Theta_{\mathcal L,h}\geq c f^*(i\,dt\wedge d\bar t)$
on a nonempty open subset of the central fiber and its neighborhood,
where $c>0$ and $v(0)$ does not vanish identically on that subset, then the
curvature pairing of $\mathcal E$ with $v(0)$ is strictly positive at $0$.
\end{lemma}

\begin{proof}
We use the representative construction in Berndtsson's proof of direct-image
positivity \cite[Proposition~4.2 and formula~(4.8)]{Berndtsson2009}; the
one-dimensional-base formulation is also
\cite[formula~(2.1) and Lemma~2.1]{Berndtsson2011}.  The calculation below
makes clear why no strict positivity along the fibers is needed.

Write $F=f^{-1}(0)$, and fix a K\"ahler form $\omega$ on $\mathcal Z$.
A smooth splitting of the relative cotangent sequence gives an absolute
lift $w$ of the relative form $v$.  Since $dt\wedge w$ is holomorphic, write
\[
 \bar\partial w=dt\wedge\eta,
 \qquad \partial^h w=dt\wedge\mu,
\]
where $\partial^h$ is the $(1,0)$ part of the Chern connection.  On $F$,
$\eta$ has type $(r-1,1)$ and $\mu$ has type $(r,0)$.
Differentiating the fiberwise pairings of $v$ with local holomorphic
sections of $\mathcal E$ identifies the projection of $\mu$ onto the
holomorphic $\mathcal L|_F$-valued top forms with $D'v(0)$.  These sections
exhaust the top forms on $F$ by base change.  Thus $\mu|_F$ is orthogonal
to the holomorphic top forms.  The terms obtained by differentiating the
conjugate relative section have a $d\bar t$ factor and do not contribute
to this $(1,0)$ derivative.

Furthermore, $(\omega\wedge\eta)|_F$ is $\bar\partial$-exact.  Indeed
$w\wedge\omega$, of type $(r+1,1)$ on the total space, is divisible by
$dt$; applying $\bar\partial$ and using $d\omega=0$ gives the assertion
on $F$.  We can therefore change the lift to $w+dt\wedge\beta$, where
$\beta$ has type $(r-1,0)$, so that
\begin{equation}\label{coa:primitive-choice}
 \mu|_F=0,\qquad (\omega\wedge\eta)|_F=0
\end{equation}
for the changed lift.  To verify simultaneous solvability, put
$\chi=\omega\wedge\beta$.  Wedge multiplication by $\omega$ is an
isomorphism from $(r-1,0)$-forms to $(r,1)$-forms on $F$, and the
K\"ahler identity gives
\[
 \bar\partial^*\chi=i\partial^h\beta.
\]
The required equations are
$\bar\partial^*\chi=i\mu$ and
$\bar\partial\chi=\omega\wedge\eta$.
The first right-hand side lies in the image of $\bar\partial^*$ because
of the preceding orthogonality; the second is $\bar\partial$-exact.
Let $G_{\bar\partial}$ denote the Green operator of the coefficient
$\bar\partial$-Laplacian on the compact fiber.  With
$q=\omega\wedge\eta$, the two forms
\[
 \chi_1=\bar\partial G_{\bar\partial}(i\mu),\qquad
 \chi_2=\bar\partial^*G_{\bar\partial}q
\]
satisfy
$\bar\partial^*\chi_1=i\mu$, $\bar\partial\chi_1=0$,
$\bar\partial\chi_2=q$, and $\bar\partial^*\chi_2=0$.
Indeed $i\mu$ is orthogonal to the harmonic $(r,0)$-forms and $q$ is
$\bar\partial$-exact.  Thus $\chi=\chi_1+\chi_2$ solves both equations.
Inverting $\omega\wedge$ and extending $\beta$ smoothly off $F$ proves
\eqref{coa:primitive-choice}.  When $r=1$, $q=0$ by type and the same
construction applies.

For this lift, differentiation of
$\|v\|^2=c_r f_*(w\wedge\overline w\,h)$ gives
\[
 \partial\|v\|^2
   =c_r f_*(\partial^h w\wedge\overline w\,h).
\]
Apply $i\bar\partial$, use
$\bar\partial\partial^h w
 =\Theta_{\mathcal L,h}\wedge w-\partial^h\bar\partial w$,
and note that the quadratic term in $\partial^h w$ vanishes on $F$.
To handle the remaining derivative, apply $\partial$ to
$f_*(\bar\partial w\wedge\overline w\,h)=0$; that pushforward vanishes
by type and its $dt$ factor.  At $0$ the resulting identity is
\begin{equation}\label{coa:primitive-identity}
 \begin{split}
 i\bar\partial\partial\|v\|^2
 &= f_*\bigl(c_r w\wedge\overline w\,h
                         \wedge i\Theta_{\mathcal L,h}\bigr)\\
 &\quad-i\,dt\wedge d\bar t\,
                  c_r\int_F\eta\wedge\overline\eta\,h.
 \end{split}
\end{equation}
Metric pairing is denoted multiplicatively by $h$ in these formulas.
For a primitive $(r-1,1)$-form the last term has nonnegative sign and is
a positive normalization of its squared norm.  This proves
\eqref{coa:curvature-bound}.

The integrand on the right of \eqref{coa:curvature-bound} is nonnegative:
a form of type $(r,0)$ in complex dimension $r+1$ is decomposable, and
wedging its squared form with a semipositive $(1,1)$-form gives a
nonnegative top-degree form.  On the stated patch we can replace the
curvature by $c f^*(i\,dt\wedge d\bar t)$ in a lower bound.  Only the
relative restriction of $w$ then contributes.  Its integral over an open
set on which $v(0)$ is nonzero is positive.  This proves strictness.
\end{proof}

For a higher-dimensional base we apply Lemma~\ref{coa:curvature} to small
embedded disks in its tangent directions.  In rank one a local holomorphic
frame may be chosen Chern-normal at a specified point.  Dividing
\eqref{coa:curvature-bound} by its squared norm then gives the curvature
of minus the logarithm of that norm.  A lower bound by the pullback of a
positive base form on one patch yields strict positivity in every
nonzero tangent direction at the image point.  If the relative dimension
is zero and the generic rank is one, the smooth map is locally an
isomorphism and this assertion follows directly from the coefficient
metric.

\subsection{A coarea estimate}

\begin{lemma}\label{coa:coarea}
Let $f:Z\to Y$ be a surjective morphism between connected smooth projective
varieties, with $\dim Z=n$ and $\dim Y=b>0$.  Let $\mu$ be a smooth
nonnegative top-degree form on $Z$.  Then $f_*\mu$ is absolutely continuous
with respect to any smooth positive volume on $Y$, and its density belongs
to $L^{1+\eta}(Y)$ for some $\eta>0$.
\end{lemma}

\begin{proof}
Fix K\"ahler metrics on $Z$ and $Y$, and use their volumes and the induced
fiber volumes.  On the smooth locus of $f$ let $J_f$ be its real coarea
Jacobian.  In orthonormal complex frames $J_f$ is the sum of the squared
absolute values of the top complex minors of $df$; equivalently, up to a
fixed normalization,
\[
 f^*\omega_Y^b\wedge\omega_Z^{n-b}=J_f\,\omega_Z^n.
\]
In finitely many coordinate charts it is comparable with the squared
norm of the corresponding holomorphic Jacobian-minor ideal.  This ideal
is not identically zero, because $f$ is dominant in characteristic zero.
A log resolution of the ideal reduces small inverse powers to products
of powers of coordinate absolute values.  Choosing their exponents
smaller than their finitely many integrability thresholds, and using
compactness, gives
\begin{equation}\label{coa:inverse-jacobian}
 \int_Z J_f^{-\eta}\,dV_Z<\infty
\end{equation}
for some $\eta>0$.

Let $U\subset Y$ be a dense Zariski open over which $f$ is smooth.  Since
$\mu\leq C\,dV_Z$, the coarea formula gives, for the density $\rho$,
\[
 \rho(y)\leq C\int_{Z_y}J_f^{-1}\,dV_{Z_y}
 \qquad (y\in U).
\]
The fiber volumes are constant on each connected smooth-base component,
by closedness of $\omega_Z^{n-b}$ and proper smooth local triviality.
They are in particular uniformly bounded.  The power-mean inequality
and coarea therefore give
\begin{equation}\label{coa:density-estimate}
 \begin{split}
 \int_U\rho^{1+\eta}\,dV_Y
 &\leq C'\int_U\int_{Z_y}J_f^{-1-\eta}\,dV_{Z_y}\,dV_Y\\
 &=C'\int_{f^{-1}(U)}J_f^{-\eta}\,dV_Z<\infty.
 \end{split}
\end{equation}
Finally, $f^{-1}(Y\setminus U)$ is a proper analytic subset of $Z$, so
has zero measure for the smooth form $\mu$.  There is consequently no
additional measure supported on $Y\setminus U$.  This proves both
absolute continuity and the claimed exponent.
\end{proof}

\subsection{Extension across base divisors}

\begin{proposition}\label{coa:volume-metric}
Let $X$ and $Z$ be connected smooth projective varieties, let
$g:Z\to X$ be a birational morphism, and let $f:Z\to Y$ be a surjective
morphism to a connected smooth projective variety of positive dimension.
Put $L=-K_X$, and assume $L$ has a smooth Hermitian metric $h$ with
semipositive curvature.  Let $E=K_{Z/X}$ be the effective Jacobian divisor.
Assume
\[
 \operatorname{rank} f_*\mathcal O_Z(E)=1.
\]
Let $\mathcal S$ be the set of prime divisors $G$ on $Y$ such that every
prime divisor of $Z$ dominating $G$ is $g$-exceptional.  Suppose there is
an open patch over the smooth base-change locus of $f$, disjoint from
$E$, on which the curvature of $g^*h$ dominates a positive multiple of
the pullback of a K\"ahler form on $Y$.
Then there is an effective integral divisor $A$ supported on
$\mathcal S$ and a singular metric on $M=-K_Y+A$, with psh weights $\psi$,
such that the metric is smooth and strictly positively curved on a
nonempty open subset of $Y$ and
\begin{equation}\label{coa:weighted-lp}
 e^{-\psi}|s_A|^2\in L^{1+\eta}_{\mathrm{loc}}
 \qquad\text{for some }\eta>0.
\end{equation}
Here $|s_A|$ means the absolute value of the coefficient of the canonical
section in a local holomorphic frame.
\end{proposition}

\begin{proof}
The identity section of $K_X+L$ and $h$ determine a smooth positive volume
form $\mu_X$ on $X$.  Write $\rho$ for the local density of
$f_*(g^*\mu_X)$ in base coordinates.  On a dense open $U\subset Y$, the
map $f$ is smooth, base change holds, and the canonical section $s_E$
spans $H^0(Z_y,\mathcal O_{Z_y}(E))$.  It gives the identification
\[
 f_*(K_{Z/Y}+g^*L)|_U\simeq -K_Y|_U.
\]
The fiberwise $L^2$ norm of the corresponding coordinate anticanonical
frame is $\rho$.  This follows directly by contracting $s_E$ with its
conjugate through $g^*h$ and integrating along the fibers; the resulting
top-degree form is $g^*\mu_X$.  Lemma~\ref{coa:curvature} proves that
$-\log\rho$ is psh, and the assumed patch proves strictness at its image.
Smooth dependence of fiber integrals gives smoothness on $U$.

It remains to control the metric at the boundary.  Let $G$ be a prime
component of $Y\setminus U$, and choose a prime $J\subset Z$ dominating
$G$.  Put $e_J=\operatorname{mult}_J(f^*G)>0$ and
$h_J=\operatorname{mult}_J E$.  Near general points of $J$ and $G$ there
are coordinates with
\[
 G=(t_1=0),\quad J=(z_1=0),\qquad
 t_1=z_1^{e_J},\quad t_i=z_i\ (2\leq i\leq b).
\]
Indeed $J\to G$ is generically submersive, and the unit in the first
coordinate can be absorbed into $z_1$.  The remaining coordinates are
fiber coordinates.  On such a patch, the pullback volume is bounded below
by a positive constant times $|z_1|^{2h_J}$ times coordinate volume.
Changing variables in the first coordinate, and integrating over a
fixed small fiber-coordinate polydisk, gives
\begin{equation}\label{coa:ramification-lower}
 \rho(t)\geq c\,|t_1|^{2((h_J+1)/e_J-1)}
\end{equation}
near general points of $G$, on $U$.  The denominator here is the squared
Jacobian $e_J^2|z_1|^{2(e_J-1)}$.

If $G\notin\mathcal S$, choose $J$ nonexceptional for $g$.  Then $h_J=0$,
and the exponent in \eqref{coa:ramification-lower} is nonpositive.
Thus $-\log\rho$ is locally bounded above there.  If
$G\in\mathcal S$, choose an integer
$a_G\geq\max\{0,(h_J+1)/e_J-1\}$.  With
$A=\sum_{G\in\mathcal S}a_GG$, the weight
\[
 \psi=-\log\rho+\log|s_A|^2
\]
is bounded above at general boundary points.  Only finitely many
$G\in\mathcal S$ occur, since each is the image of one of the finitely
many $g$-exceptional primes.  On $U$ the added term is pluriharmonic.
The psh extension theorem therefore extends the weight across all
codimension-one points under consideration.  The remaining analytic
subset has codimension at least two, across which psh extension is
again available.  One can see the needed local upper bound by using
small disks through nearby points whose boundaries remain in a fixed
compact subset off that analytic set, and applying the maximum
principle.  The frame transformations identify these extended weights
as a metric of $-K_Y+A$.

The strict patch can be chosen away from $A$, so strict positivity is
preserved there.  Lastly, $g^*\mu_X$ is a smooth nonnegative volume form
on the compact variety $Z$.  Lemma~\ref{coa:coarea} gives an exponent
$1+\eta>1$ for its pushforward density.  Since
$e^{-\psi}|s_A|^2=\rho$ almost everywhere in the compatible frames,
this is exactly \eqref{coa:weighted-lp}.
\end{proof}

\begin{corollary}[Smooth invariant direct image]\label{coa:invariant-volume-metric}
The conclusion of Proposition~\ref{coa:volume-metric} remains valid if
its ordinary rank-one hypothesis is replaced by
\[
 \operatorname{rank}\bigl(f_*\mathcal O_Z(K_{Z/X})\bigr)^T=1,
\]
provided that a compact torus $T$ acts holomorphically on $X$ and $Z$,
acts trivially on $Y$, preserves $f$, $g$, and the smooth metric $h$,
and all canonical bundles carry their natural linearizations.
\end{corollary}

\begin{proof}
On the smooth base-change locus, averaging over normalized Haar measure
is a holomorphic idempotent on the adjoint direct image.  Invariance of
the fiber metric makes this projection orthogonal.  Its image is the
invariant line, and its kernel is the orthogonal holomorphic complement.
Consequently the line has the quotient metric of the semipositive
smooth direct image.  For a section in this line, a Chern-normal
holomorphic extension within the line is also Chern-normal in the full
bundle, so Lemma~\ref{coa:curvature} proves the same strictness statement.
The naturally linearized Jacobian section is invariant and generates
this line.  Its norm is still the density $\rho$.

For relative dimension zero the full direct image is locally a direct
sum of the coefficient lines on the sheets of an unramified cover.
Its curvature is their direct sum; the same orthogonal projection
argument applies, and the nonzero Jacobian section detects the strict
patch.  The ramification estimate, divisor correction, and coarea
calculation in Proposition~\ref{coa:volume-metric} do not use ordinary
rank one.  They therefore prove the stated extension and integrability.
\end{proof}

\subsection{From local strictness to a global strict metric}

The metric just constructed is strictly positive on one open set.
The next step uses the extra integrability exponent to introduce a
global lower curvature bound while preserving the distinguished
integrable section.

\begin{lemma}\label{coa:weighted-effectivity}
Let $Y$ be smooth connected projective of positive dimension, let
$A\geq0$ be integral, and put $M=-K_Y+A$.  Suppose $M$ has psh metric
weights $\psi$, smooth and strictly positively curved on a nonempty
open set, and
\[
 |s_A|^2e^{-\psi}\in L^{1+\eta}_{\mathrm{loc}}
 \quad\text{for some }\eta>0.
\]
Then $M$ has a singular metric with a global strict curvature lower
bound and with $|s_A|^2$ times its squared frame norm locally
integrable.  For every pseudoeffective rational divisor $W$, an
effective rational divisor $C_A$ supported on $A$ makes $W+C_A$
rationally linearly equivalent to an effective rational divisor.
\end{lemma}

\begin{proof}
The curvature current of $\psi$ has positive absolutely continuous
top-degree mass, since it is smooth and positive definite on the
specified open set.  Boucksom's volume criterion
\cite{Boucksom2002} makes $M$ big.  Write
$M\sim_{\mathbb Q}H+G$ with $H$ ample rational and $G\geq0$ rational.
A smooth positive metric on $H$ and the divisor metric of $G$ give a
metric on $M$ with weights $\chi$, curvature bounded below by a
positive multiple of a K\"ahler form, and divisorial analytic
singularities.  Log resolution and compactness give a number
$\tau>0$ such that $|s_A|^2e^{-\chi}\in L^\tau_{\mathrm{loc}}$.
Choose $\sigma\in(0,1)$ sufficiently small that
\[
 \frac{1-\sigma}{1+\eta}+\frac{\sigma}{\tau}<1.
\]
The weights $\psi_\sigma=(1-\sigma)\psi+\sigma\chi$ then have a
global strict curvature lower bound.  Moreover,
\[
 |s_A|^2e^{-\psi_\sigma}
 =\bigl(|s_A|^2e^{-\psi}\bigr)^{1-\sigma}
  \bigl(|s_A|^2e^{-\chi}\bigr)^\sigma
\]
is locally integrable by H\"older's inequality.  Equivalently,
$\mathcal J(\psi_\sigma)\supset\mathcal O_Y(-A)$.
The controlled-boundary criterion, Theorem~\ref{thm:controlled-boundary},
now gives the assertion about $W$ and $C_A$.
\end{proof}

\subsection{The smooth metric on the maximal dominated base}
\label{coa:maximal-base}\label{coa:dominated-conversion}

We now apply the smooth coarea construction to the maximal base of
Section~\ref{conv:all-dominated}.  That base already supplies the
interpolation and the pseudoeffective divisor of vertical minima.
The additional points here are a valuation proof of the model
condition and a direct intersection argument for horizontal strictness.
These allow us to use the original smooth metric, and hence retain
the exponent $1+\eta$ in Proposition~\ref{coa:volume-metric}.

Retain the hypotheses of Theorem~\ref{conv:fixed-error}, with effective
integral divisors $N_j\sim a_jL-D$, where $D$ is pseudoeffective and
$a_j$ is a strictly increasing unbounded sequence of positive integers.
A rational map is \emph{dominated by $L$} if, on a smooth resolution
$X\xleftarrow{\ g\ }Z\xrightarrow{\ f\ }Y$, one has
\[
 g^*L-\epsilon f^*H\quad\hbox{pseudoeffective}
\]
for an ample $H$ on the projective base and a rational $\epsilon>0$.
Include the point map.  This is the map formulation of the admissible
systems used in Section~\ref{conv:all-dominated}.  Indeed, an
admissible system gives this inequality after resolving its base
ideal.  Conversely, after multiplying $H$ so that it defines a
projective embedding, its pulled-back sections push down to a
subsystem of $|g_*f^*H|$; pushing down the domination inequality
makes that system admissible.

The condition persists on higher source resolutions and under
birational base modifications.  It also persists under a generically
finite change of base to which the rational map lifts dominantly:
the pullback of the old ample class is big and dominates a small
ample class.  Finally, the joint image of two dominated maps is
dominated.  On a common source resolution, add suitable positive
multiples of their inequalities and use the sum of the pulled-back
polarizations, whose restriction to the joint image is ample.
Thus the maximality construction in Section~\ref{conv:all-dominated}
provides a base field $\mathbb C(Y)$, relatively algebraically closed
in $\mathbb C(X)$, containing the functions supplied by every
dominated map.

For the resulting smooth projective diagram put
\[
 P=\frac{g^*N_2-g^*N_1}{a_2-a_1}\sim_{\mathbb Q}g^*L,
 \qquad B=a_1P-g^*N_1.
\]
The interpolation argument of Section~\ref{conv:all-dominated} gives
$\mathbb Q$-principal divisors $U_j$ on $Y$ with
\begin{equation}\label{coa:interpolation}
 g^*N_j=a_jP-B+f^*U_j,\qquad U_1=U_2=0.
\end{equation}
For $j>2$, the two effective divisors in
$(a_j-a_1)N_2\sim(a_j-a_2)N_1+(a_2-a_1)N_j$ lie in a
multiple of the admissible system $N_2$, so their ratio belongs to
$\mathbb C(Y)$.  In particular $P$ has nonnegative horizontal
part.  If $Y$ is a point, pushing $P$ down already proves the theorem.
Assume henceforth that $b=\dim Y>0$.

\subsubsection*{Preparing vertical primes by their valuations}

The model can be arranged so that every vertical prime not exceptional
over $X$ dominates a base prime.  The flattening argument of
Section~\ref{conv:section} gives one construction.  The following
valuation argument gives another and explains why a base modification
suffices.

On an initial model only finitely many source primes map into base
codimension at least two: take an open set over which $f$ is
equidimensional and inspect the divisorial components of its
complement.  The valuation of such a prime restricts nontrivially to
$\mathbb C(Y)$, since a local function vanishing on its proper center
has positive order.  The restricted discrete valuation is divisorial.
To see this, its residue transcendence degree can drop by at most
$\dim X-\dim Y$.  Residues algebraically independent over the
restricted residue field lift to elements algebraically independent
over $\mathbb C(Y)$: a polynomial relation, with coefficients scaled
to have minimum valuation zero, would give a residue relation.
The source residue field has transcendence degree $\dim X-1$.
The restricted residue field therefore has transcendence degree at
least $b-1$, and the valuation inequality gives the opposite bound.

Choose a projective base model on which these finitely many valuations
have divisorial centers, and resolve it.  One can construct such a
model by resolving maps to copies of $\mathbb P^1$ defined by
$b-1$ independent residues: the center then has dimension at least
$b-1$ and remains proper.  Take a common smooth resolution of the
previous source and the graph over the new base.  Primes already
dominating base divisors retain that property.  New source primes
exceptional over the old source are also exceptional over $X$.
This proves the required preparation.  The domination and field
properties persist by the preceding observations.  Retain $P,B$ for
their pullbacks to this model and $U_j$ for the corresponding principal
divisors on the modified base; \eqref{coa:interpolation} is unchanged.

On this prepared model define
\begin{equation}\label{coa:minima}
 w_G=\min_{J\mapsto G}\frac{\operatorname{mult}_JP}
                              {\operatorname{mult}_Jf^*G},
 \qquad W=\sum_Gw_GG.
\end{equation}
The sum has finite support.  Formula~\eqref{coa:interpolation} makes
$W$ pseudoeffective by the same fixed-error limit as in
Section~\ref{conv:all-dominated}: if
$c_G=\min_{J\mapsto G}\operatorname{mult}_JB/
\operatorname{mult}_Jf^*G$ and $C=\sum_Gc_GG$, then
$a_jW+U_j-C\ge0$; dividing by $a_j$ gives numerical classes
converging to $[W]$.

Let $\mathcal S$ consist of base primes whose dominating source
primes are all $g$-exceptional.  It is finite, and the model condition
makes the entire pullback of each member exceptional over $X$.
Lemma~\ref{conv:return}, with its $R=P$ and $M=W$, reduces the
proof to making $W+C_{\mathcal S}$ rationally effective for an
effective correction supported on $\mathcal S$.

\subsubsection*{Rank one and horizontal strictness}

Put $E=K_{Z/X}$.  The maximality argument in
Section~\ref{conv:all-dominated} gives
\begin{equation}\label{coa:rank-one}
 \operatorname{rank}f_*\mathcal O_Z(E)=1.
\end{equation}
Its concrete input is a domination inequality
$b'g^*L-f^*H'$ pseudoeffective for a sufficiently ample base
polarization $H'$.  Two generic-fiber independent adjoint sections
would, after an ample base twist, push down to an admissible pencil
with ratio outside $\mathbb C(Y)$, contradicting the defining
maximality.  The Jacobian section supplies the nonzero rank-one
generator.

Let $\alpha\ge0$ be the smooth curvature of the original pulled-back
metric on $g^*L$, normalized to represent $g^*c_1(L)$, and let $r$
be its maximum rank.  If $r=n=\dim Z$, then $L$ is nef with positive
top self-intersection and hence big, giving the required nonvanishing.
Suppose $r<n$.  Choose K\"ahler
forms $\omega_Z,\omega_Y$ representing ample classes.
Domination gives $b'[\alpha]-[f^*\omega_Y]$ pseudoeffective for
some $b'>0$.  Intersecting with
$[\alpha]^r[\omega_Z]^{n-r-1}$, and using $\alpha^{r+1}=0$, yields
\[
 \int_Z f^*\omega_Y\wedge\alpha^r\wedge\omega_Z^{n-r-1}=0.
\]
The integrand is nonnegative, so it vanishes pointwise.  At a
rank-$r$ point, this says that $f^*\omega_Y$ vanishes on
$\ker\alpha$.  A nonnegative Hermitian form has zero mixed pairings
against any vector on which its quadratic value vanishes.  Hence
\begin{equation}\label{coa:kernel-inclusion}
 \ker\alpha\subseteq\ker df.
\end{equation}
At a smooth point this forces $r\ge b$; rank zero is therefore
excluded.  The maximum-rank open set meets both the complement of
$E$ and the smooth base-change locus.  On a smaller relatively compact
constant-rank patch there, \eqref{coa:kernel-inclusion} gives
$\alpha\ge c f^*\omega_Y$ for some $c>0$.

Proposition~\ref{coa:volume-metric} now applies with its birational
map $g:Z\to X$ and fibration $f:Z\to Y$.  It supplies an integral
effective $A$ supported on $\mathcal S$ and a semipositive metric
on $-K_Y+A$, smooth and strictly positive on a nonempty open set,
whose weighted density belongs to $L^{1+\eta}_{\rm loc}$.
Lemma~\ref{coa:weighted-effectivity} gives an effective $C_A$
supported on $A$ with $W+C_A$ rationally effective.
The return lemma applies to $P\sim_{\mathbb Q}g^*L$: its horizontal
part is nonnegative, its minimum divisor is $W$, and every component
of $f^*C_A$ is exceptional over $X$.  It gives the desired section
of a positive multiple of $L$.

This proves Theorem~\ref{conv:fixed-error} using the original smooth
coefficient metric.  The determinant and tensor consequences remain
Lemma~\ref{conv:determinant} and Corollary~\ref{conv:tensor-conversion}.
The next section retains a different bigness proof and a
fiber-intersection test for horizontal positivity.
